% Figura extraída de "Tema 2.tex" (línea 540 del original)
\documentclass[12pt,border=1pt]{standalone}
\usepackage[spanish,es-tabla]{babel}
\usepackage[utf8]{inputenc}
\usepackage{amsmath}
\usepackage{xcolor}
\usepackage{tikz}
\usepackage{pgfplots}
\usepackage{circuitikz}
\usetikzlibrary{babel, circuits.ee}

\definecolor{applegreen}{rgb}{0.55, 0.71, 0.0}
\definecolor{airforceblue}{rgb}{0.36, 0.54, 0.66}
\definecolor{amber}{rgb}{1.0, 0.49, 0.0}

\newcommand{\diff}[3][1]{%
	\IfBlankTF{#1}{\frac{\mathrm{d}{#2}}{\mathrm{d}{#3}}}%
	{\frac{\mathrm{d}^{#1} {#2}}{\mathrm{d}{#3}^{#1}}}%
}%
\newcommand{\pdiff}[3][1]{%
	\IfBlankTF{#1}{\frac{\partial {#2}}{\partial {#3}}}%
	{\frac{\partial^{#1} {#2}}{\partial {#3}^{#1}}}%
}%

\begin{document}
\begin{tikzpicture}
%% DIBUJAMOS EL CUBO
% Dimensiones
\def\L{8} % Longitud (x)
\def\W{2} % Altura (y)
\def\H{3} % Ancho (z)

% Dibuja los lados exteriores
\draw  (0,\W,0) -- (0,\W,\H); %Lado izquierdo arriba
\draw (0,0,\H)--(0,\W,\H); % Lado izquierdo frontal
\draw (\L,0,0) -- (\L,\W,0) -- (\L,\W,\H); %Lado derecho trasero y superior derecho
\draw (0,0,\H) -- (\L,0,\H) -- (\L,\W,\H); %Lado bajo frontal y derecho frontal
\draw (0,\W,0) -- (\L,\W,0); %Lado de arriba atrás
\draw (0,\W,\H)-- (\L,\W,\H); % Lado de abajo adelante
\draw (\L,0,0) -- (\L,0,\H); %Lado derecho abajo
% Dibuja los lados interiores punteados
\draw[dashed] (0,0,0) -- (0,0,\H);
\draw[dashed] (0,0,0) -- (\L,0,0);
\draw[dashed] (0,0,0)-- (0,\W,0);
%%
%%% Ponemos las dimensiones
    %Longitud
\draw[latex-latex] (0, 0, \H+1) -- (\L, 0, \H+1) node[midway, below, yshift=-5] {$L$};
\draw (0,0,\H+0.5) --(0,0,\H+1.5);
\draw (\L,0,\H+0.5) --(\L,0,\H+1.5);
    %Anchura
\draw[latex-latex] (\L+0.5, 0, 0) -- (\L+0.5, 0, \H);
\node at (\L+1,-0.1,\H/2-0.1) {$W$};
\draw (\L+0.3,0,\H) --(\L+0.75,0,\H);
\draw (\L+0.3,0,0) --(\L+0.75,0,0);
    %Altura
\draw[latex-latex] (\L+0.5, 0, 0) -- (\L+0.5, \W,0);
\node at (\L+1,\W/2,0) {$H$};
\draw (\L+0.3,\W,0) --(\L+0.75,\W,0);
%%%% Corrientes campos etc
\draw [-latex, blue, thick] (7,\W/3,\H/2)--(0,\W/3,\H/2);
\node at (0.7, \W/3-0.5, \H/2-0.5) {\textcolor{blue}{$\mathbf{J}$}};
\draw [-latex, thick, red] (\L/2,\W/3,\H/2)--(\L/2,\W+1.3,\H/2) ;
\node at (\L/2,\W+0.5,0) {\textcolor{red}{$\mathbf{B_z}$}};
\draw[-latex, thick] (\L/2,\W/3,\H/2)--(\L/2,\W/3,\H+1.3);
\node at (\L/2+1,\W/3,\H+0.4) {$\mathbf{B_z}\times \mathbf{J} $};
\draw [-latex] (1.2, \W-0.3, \H-0.1)--(2, \W-0.3, \H-0.1);
\node at (2.6, \W-0.3, \H-0.1) {$\mathbf{v_x}$};
  % Dibuja los ejes
  \draw[-latex, ultra thick] (-4,0,0) -- (-5,0,0) node[below]{$X$};
  \draw[-latex, ultra thick] (-4,0,0) -- (-4,1,0) node[right]{$Z$};
  \draw[-latex, ultra thick] (-4,0,0) -- (-4,0,1) node[right]{$ \  \ Y$};
\end{tikzpicture}
\end{document}
